\documentclass[12pt,twoside,openright]{report}

% Packages
\usepackage[utf8]{inputenc}
\usepackage[T1]{fontenc}
\usepackage{amsmath,amssymb,amsthm}
\usepackage{mathtools}
\usepackage{algorithm}
\usepackage{algpseudocode}
\usepackage{graphicx}
\usepackage{subcaption}
\usepackage{listings}
\usepackage{xcolor}
\usepackage{hyperref}
\usepackage{cleveref}
\usepackage{tikz}
\usepackage{pgfplots}
\pgfplotsset{compat=1.18}
\usepackage{geometry}
\usepackage{fancyhdr}
\usepackage{natbib}
\usepackage{nomencl}
\usepackage{booktabs}
\usepackage{multirow}
\usepackage{tocloft}

% Page layout
\geometry{
    a4paper,
    left=1.5in,
    right=1in,
    top=1in,
    bottom=1in,
}

% Header and footer
\pagestyle{fancy}
\fancyhf{}
\fancyhead[LE]{\leftmark}
\fancyhead[RO]{\rightmark}
\fancyfoot[C]{\thepage}
\renewcommand{\headrulewidth}{0.4pt}

% Theorem environments
\theoremstyle{plain}
\newtheorem{theorem}{Theorem}[chapter]
\newtheorem{lemma}[theorem]{Lemma}
\newtheorem{proposition}[theorem]{Proposition}
\newtheorem{corollary}[theorem]{Corollary}

\theoremstyle{definition}
\newtheorem{definition}[theorem]{Definition}
\newtheorem{example}[theorem]{Example}
\newtheorem{remark}[theorem]{Remark}
\newtheorem{assumption}[theorem]{Assumption}

% Code listing style
\lstdefinestyle{pythonstyle}{
    language=Python,
    basicstyle=\ttfamily\small,
    keywordstyle=\color{blue}\bfseries,
    commentstyle=\color{gray}\itshape,
    stringstyle=\color{red},
    numbers=left,
    numberstyle=\tiny\color{gray},
    stepnumber=1,
    numbersep=5pt,
    backgroundcolor=\color{white},
    showspaces=false,
    showstringspaces=false,
    showtabs=false,
    frame=single,
    rulecolor=\color{black},
    tabsize=4,
    captionpos=b,
    breaklines=true,
    breakatwhitespace=false,
    escapeinside={\%*}{*)},
}

\lstset{style=pythonstyle}

% Nomenclature
\makenomenclature

% Hyperref setup
\hypersetup{
    colorlinks=true,
    linkcolor=blue,
    citecolor=blue,
    filecolor=magenta,
    urlcolor=cyan,
    pdftitle={Adaptive Control with Memory Kernels: Lyapunov Stability and Convergence Analysis},
    pdfauthor={Your Name},
}

% Custom commands
\newcommand{\R}{\mathbb{R}}
\newcommand{\C}{\mathbb{C}}
\newcommand{\N}{\mathbb{N}}
\newcommand{\norm}[1]{\left\| #1 \right\|}
\newcommand{\abs}[1]{\left| #1 \right|}
\newcommand{\inner}[2]{\left\langle #1, #2 \right\rangle}
\newcommand{\deriv}[2]{\frac{d #1}{d #2}}
\newcommand{\pderiv}[2]{\frac{\partial #1}{\partial #2}}
\newcommand{\Lp}[1]{L^{#1}}
\newcommand{\thetahat}{\hat{\theta}}
\newcommand{\thetastar}{\theta^*}
\newcommand{\thetatilde}{\tilde{\theta}}

% Title information
\title{
    {\Large Adaptive Control with Memory Kernels:}\\[0.5cm]
    {\huge Lyapunov Stability and Convergence Analysis}\\[0.5cm]
    {\large with Applications to Physiological Systems}
}
\author{Your Name}
\date{\today}

\begin{document}

% Front matter
\frontmatter
\pagestyle{empty}

% Title page
\begin{titlepage}
    \centering
    \vspace*{2cm}

    {\LARGE\bfseries Adaptive Control with Memory Kernels:\par}
    \vspace{0.5cm}
    {\Huge\bfseries Lyapunov Stability and Convergence Analysis\par}
    \vspace{0.5cm}
    {\Large with Applications to Physiological Systems\par}

    \vspace{2cm}

    {\Large A Thesis Presented\par}
    \vspace{0.3cm}
    {\Large by\par}
    \vspace{0.5cm}
    {\LARGE\bfseries Your Name\par}

    \vspace{2cm}

    {\large to\par}
    \vspace{0.3cm}
    {\large The Department of Control Systems and Bioengineering\par}

    \vspace{1cm}

    {\large in partial fulfillment of the requirements\par}
    {\large for the degree of\par}
    \vspace{0.3cm}
    {\Large\bfseries Doctor of Philosophy\par}

    \vspace{1cm}

    {\large in the subject of\par}
    \vspace{0.3cm}
    {\Large Control Theory and Adaptive Systems\par}

    \vfill

    {\large Your University\par}
    {\large City, Country\par}
    \vspace{0.3cm}
    {\large \today\par}

\end{titlepage}


% Abstract
\chapter*{Abstract}
\addcontentsline{toc}{chapter}{Abstract}

This thesis presents a comprehensive treatment of adaptive control systems incorporating memory kernels, with rigorous Lyapunov-based stability and convergence analysis. We develop a unified theoretical framework for systems with distributed parameter uncertainty, delayed coupling, and non-Markovian dynamics.

The primary contributions of this work include:

\begin{enumerate}
    \item \textbf{Mathematical Framework:} A complete functional-analytic foundation for adaptive control with memory, including detailed treatment of Barbalat's Lemma, persistent excitation (PE) theory, and Sobolev function spaces.

    \item \textbf{Lyapunov Stability Theory:} Novel Lyapunov-Krasovskii functionals that incorporate memory kernel effects, with full algebraic derivations demonstrating global asymptotic stability (GAS) under PE conditions and uniform ultimate boundedness (UUB) with bounded disturbances.

    \item \textbf{Adaptive Update Laws:} Formal derivation of projection-based and saturation-based adaptive parameter update laws, with convergence guarantees and robustness to measurement noise and unmatched uncertainties.

    \item \textbf{Robust Extensions:} Systematic treatment of practical considerations including non-PE scenarios, state-dependent disturbances, and unmodeled dynamics, with explicit bounds on tracking error and parameter estimation error.

    \item \textbf{Physiological Applications:} Application to heart-brain coupling models and neural control systems, demonstrating the practical utility of the theoretical framework in biomedical engineering contexts.

    \item \textbf{Computational Implementation:} Complete Python implementation with numerical validation, including comparison of PE versus non-PE scenarios and demonstration of theoretical convergence rates.
\end{enumerate}

The thesis begins with a comprehensive literature review covering classical adaptive control (Narendra, Åström), Lyapunov stability theory (Khalil, Slotine), and memory kernel modeling. Chapter 3 develops the mathematical preliminaries, proving Barbalat's Lemma from first principles and establishing the function space framework.

Chapters 4-7 form the theoretical core, defining the system class, constructing augmented memory states, deriving the complete Lyapunov candidate functional, and proving the main stability theorems. Every algebraic step is shown explicitly to enable verification and extension by future researchers.

Chapter 8 addresses robust extensions critical for practical implementation, including adaptive safety mechanisms for explainable AI (xAI) applications. Chapter 9 presents extensive simulation results with Python/MATLAB code listings, demonstrating parameter convergence, boundedness properties, and comparative performance under PE and non-PE conditions.

The thesis concludes with practical implementation checklists for convergence verification, robustness assessment, and safety certification, making the theoretical results accessible to practitioners in control engineering, robotics, and biomedical systems.

All proofs are constructive and complete, with detailed appendices providing full derivations, additional lemmas, and commented source code for all numerical experiments.

\vspace{1cm}

\noindent\textbf{Keywords:} Adaptive control, Lyapunov stability, memory kernels, persistent excitation, Barbalat's Lemma, uniform ultimate boundedness, physiological systems, heart-brain coupling


% Dedication (optional)
% \include{chapters/dedication}

% Acknowledgments (optional)
% \include{chapters/acknowledgments}

\pagestyle{fancy}

% Table of Contents
\tableofcontents
\addcontentsline{toc}{chapter}{Contents}

% List of Figures
\listoffigures
\addcontentsline{toc}{chapter}{List of Figures}

% List of Tables
\listoftables
\addcontentsline{toc}{chapter}{List of Tables}

% Nomenclature
\printnomenclature
\addcontentsline{toc}{chapter}{Nomenclature}

% Main matter
\mainmatter

% Chapters
\chapter{Introduction}
\label{ch:introduction}

\section{Motivation}

Adaptive control addresses the fundamental challenge of controlling systems with uncertain or time-varying parameters. Unlike robust control, which attempts to handle uncertainty through conservative design margins, adaptive control \textit{learns} the unknown parameters online, potentially achieving superior performance while maintaining stability guarantees.

The need for adaptive control arises in numerous applications:
\begin{itemize}
    \item \textbf{Aerospace:} Aircraft dynamics change with fuel consumption, payload, and atmospheric conditions.
    \item \textbf{Robotics:} Manipulators handle objects with unknown mass and inertia.
    \item \textbf{Biomedical Systems:} Patient-specific physiological parameters vary widely and evolve over time.
    \item \textbf{Automotive:} Vehicle dynamics depend on load, tire wear, and road conditions.
    \item \textbf{Process Control:} Chemical plants exhibit parameter drift due to catalyst aging and feedstock variations.
\end{itemize}

\subsection{The Role of Memory in Adaptive Systems}

Classical adaptive control theory, developed primarily in the 1970s--1990s by pioneers such as Narendra, Åström, and Slotine, focused on memoryless (Markovian) systems where the current state fully determines future evolution. However, many physical systems exhibit \textbf{memory effects}:
\begin{itemize}
    \item \textbf{Biological Systems:} Neural plasticity, muscle fatigue, and cardiovascular dynamics involve distributed time delays and history-dependent responses.
    \item \textbf{Viscoelastic Materials:} Stress-strain relationships depend on deformation history through memory kernels.
    \item \textbf{Financial Systems:} Market dynamics incorporate memory of past prices and trading volumes.
    \item \textbf{Climate Models:} Atmospheric and oceanic processes exhibit long-range temporal correlations.
\end{itemize}

Memory effects are mathematically captured by integro-differential equations of the form:
\begin{equation}
    \dot{x}(t) = f(x(t), u(t)) + \int_{-\tau}^0 K(\xi) g(x(t+\xi)) \, d\xi,
\end{equation}
where the kernel $K(\cdot)$ weights past states $x(t+\xi)$ for $\xi \in [-\tau, 0]$.

\textbf{Challenge:} Standard Lyapunov stability theory does not directly apply to systems with memory. We require \textit{Lyapunov-Krasovskii} functionals that account for the entire state history over $[t-\tau, t]$.

\subsection{Contributions of This Thesis}

This thesis makes the following contributions to adaptive control theory:

\begin{enumerate}
    \item \textbf{Unified Mathematical Framework:} We develop a complete functional-analytic foundation incorporating:
    \begin{itemize}
        \item Sobolev spaces $W^{k,p}$ for state regularity,
        \item Barbalat's Lemma with full proof from first principles,
        \item Persistent excitation theory with spectral characterization.
    \end{itemize}

    \item \textbf{Lyapunov-Krasovskii Stability Analysis:} We construct novel Lyapunov functionals for adaptive systems with memory, providing:
    \begin{itemize}
        \item Complete algebraic derivation of time derivatives (every step shown),
        \item Proof of uniform ultimate boundedness (UUB) under bounded disturbances,
        \item Proof of global asymptotic stability (GAS) under persistent excitation.
    \end{itemize}

    \item \textbf{Formal Adaptive Update Laws:} We derive projection-based and saturation-based parameter update laws with:
    \begin{itemize}
        \item Convergence guarantees,
        \item Robustness to measurement noise,
        \item Computational efficiency.
    \end{itemize}

    \item \textbf{Robust Extensions for Practical Implementation:} We address:
    \begin{itemize}
        \item Non-PE scenarios (tracking without parameter convergence),
        \item State-dependent disturbances,
        \item Unmodeled dynamics,
        \item Adaptive safety mechanisms for explainable AI (xAI) applications.
    \end{itemize}

    \item \textbf{Physiological Application:} We apply the framework to \textit{heart-brain coupling models}, demonstrating:
    \begin{itemize}
        \item Adaptive control of cardiac rhythm via neural stimulation,
        \item Brain-computer interface (BCI) integration with OpenBCI and Neuralink-style hardware,
        \item Parameter identification in physiological oscillators (FitzHugh-Nagumo, Van der Pol).
    \end{itemize}

    \item \textbf{Complete Computational Implementation:} We provide:
    \begin{itemize}
        \item Python reference implementation with full documentation,
        \item Simulation results validating theoretical predictions,
        \item Comparison of PE versus non-PE scenarios,
        \item Practical implementation checklists.
    \end{itemize}
\end{enumerate}

\section{Historical Context}

\subsection{Classical Adaptive Control (1960s--1980s)}

The foundations of adaptive control were laid in the 1960s--1970s:
\begin{itemize}
    \item \textbf{Model Reference Adaptive Control (MRAC):} Whitaker (1958), Parks (1966), and Monopoli (1974) developed MRAC for systems tracking a reference model.
    \item \textbf{Self-Tuning Regulators:} Åström and Wittenmark (1973) introduced self-tuning regulators based on recursive parameter estimation.
    \item \textbf{Stability Theory:} Narendra and colleagues established Lyapunov-based stability proofs in the 1970s--1980s \cite{Narendra1989}.
\end{itemize}

\textbf{Key Insight:} The certainty equivalence principle—use parameter estimates $\thetahat(t)$ as if they were true parameters—can be stabilized through appropriately designed adaptive laws.

\subsection{Persistent Excitation (1970s--1980s)}

A crucial discovery was that tracking error convergence ($e \to 0$) does not guarantee parameter convergence ($\thetahat \to \thetastar$). Parameter convergence requires \textbf{persistent excitation}:
\begin{itemize}
    \item Anderson (1977) formalized PE conditions.
    \item Narendra and Annaswamy (1987) provided detailed PE analysis for MRAC.
    \item Boyd and Sastry (1986) gave spectral characterization of PE signals.
\end{itemize}

\subsection{Robust Adaptive Control (1980s--1990s)}

Practical implementation revealed challenges:
\begin{itemize}
    \item \textbf{Unmodeled Dynamics:} Real systems deviate from nominal models.
    \item \textbf{Disturbances:} External perturbations and measurement noise.
    \item \textbf{Computational Constraints:} Real-time parameter adaptation.
\end{itemize}

Robust adaptive control, developed by Ioannou and Sun (1996), Slotine and Li (1991), and Krstić, Kanellakopoulos, and Kokotović (1995), introduced:
\begin{itemize}
    \item $\sigma$-modification and $e$-modification to prevent parameter drift,
    \item Projection algorithms to enforce parameter bounds,
    \item Deadzone modifications for robustness to disturbances.
\end{itemize}

\subsection{Adaptive Control with Memory (1990s--Present)}

Interest in non-Markovian systems led to:
\begin{itemize}
    \item \textbf{Fractional-Order Systems:} Fractional calculus captures memory through non-integer derivatives.
    \item \textbf{Time-Delay Systems:} Explicit time delays in feedback (Niculescu 2001, Gu et al. 2003).
    \item \textbf{Integro-Differential Systems:} Memory kernels in viscoelasticity, population dynamics, and epidemiology.
\end{itemize}

\textbf{Gap in Literature:} Most work treats memory as a perturbation or focuses on specific kernel forms (exponential, polynomial). A \textit{unified framework} for general memory kernels with adaptive parameter uncertainty is lacking.

\textbf{This thesis fills this gap.}

\section{Physiological Application: Heart-Brain Coupling}

A motivating application is the \textbf{heart-brain coupling model (HBCM)}, which describes bidirectional interactions between cardiac and neural oscillators:
\begin{align}
    \text{Neural (FitzHugh-Nagumo):} \quad &
    \begin{cases}
        \dot{v} = v - v^3/3 - w + I_{\text{cardiac} \to \text{neural}} \\
        \dot{w} = (v + a - bw) / c
    \end{cases} \\[0.5em]
    \text{Cardiac (Van der Pol):} \quad &
    \begin{cases}
        \dot{x} = y \\
        \dot{y} = -\omega^2 x + \mu (1 - x^2) y + I_{\text{neural} \to \text{cardiac}}
    \end{cases}
\end{align}

The coupling terms involve delayed feedback:
\begin{align}
    I_{\text{cardiac} \to \text{neural}}(t) &= g_1 x(t - \tau_1), \\
    I_{\text{neural} \to \text{cardiac}}(t) &= g_2 v(t - \tau_2),
\end{align}
where $g_1, g_2$ are coupling gains and $\tau_1, \tau_2$ are propagation delays (typically 50--200 ms).

\subsection{Control Objective}

\textbf{Goal:} Regulate cardiac rhythm (heart rate, heart rate variability) via neural stimulation, adapting to unknown patient-specific parameters $(a, b, c, \mu, \omega, g_1, g_2, \tau_1, \tau_2)$.

\textbf{Challenges:}
\begin{itemize}
    \item \textbf{Parameter Uncertainty:} Physiological parameters vary across individuals and over time.
    \item \textbf{Delays:} Axonal conduction and synaptic delays introduce memory effects.
    \item \textbf{Safety:} Neural stimulation must remain within safe bounds.
    \item \textbf{Real-Time Adaptation:} Parameters must be estimated from online measurements (EEG, ECG).
\end{itemize}

This thesis develops adaptive controllers for HBCM, integrating:
\begin{itemize}
    \item Brain-computer interfaces (OpenBCI, Neuralink) for neural state measurement,
    \item Adaptive Lyapunov-based control for cardiac regulation,
    \item Safety constraints via projection and saturation.
\end{itemize}

\section{Thesis Organization}

The remainder of this thesis is organized as follows:

\begin{description}
    \item[Chapter 2: Literature Review] Comprehensive survey of adaptive control, Lyapunov stability theory, memory kernel modeling, and physiological applications.

    \item[Chapter 3: Mathematical Preliminaries] Function spaces ($\Lp{p}$, $W^{k,p}$), Barbalat's Lemma with complete proof, persistent excitation theory.

    \item[Chapter 4: System Definition and Augmented States] Formal problem statement, uncertainty modeling, augmented memory state construction.

    \item[Chapter 5: Lyapunov-Krasovskii Stability Analysis] Construction of Lyapunov functional, complete derivative algebra, stability conditions.

    \item[Chapter 6: Adaptive Update Laws] Derivation of gradient-based, projection-based, and saturation-based adaptive laws with convergence analysis.

    \item[Chapter 7: Main Stability Theorems] Formal statements and proofs of GAS under PE, UUB with disturbances, convergence rates.

    \item[Chapter 8: Robust Extensions] Non-PE scenarios, unmodeled dynamics, disturbance rejection, adaptive safety for xAI.

    \item[Chapter 9: Simulation Results] Python/MATLAB implementations, numerical experiments, PE vs. non-PE comparisons, parameter convergence plots.

    \item[Chapter 10: Practical Guidelines] Implementation checklists, tuning heuristics, convergence diagnostics, safety verification.

    \item[Chapter 11: Conclusion] Summary of contributions, future research directions.

    \item[Appendices] Detailed derivations, additional lemmas, complete code listings, extended results.
\end{description}

\section{Notation and Conventions}

Throughout this thesis, we use the following notation:
\begin{itemize}
    \item $\R, \C, \N$: Real numbers, complex numbers, natural numbers
    \item $\R^n$: $n$-dimensional Euclidean space
    \item $\norm{\cdot}$: Euclidean norm for vectors, induced 2-norm for matrices
    \item $\Lp{p}[0,\infty)$: Lebesgue space with norm $\norm{f}_{\Lp{p}} = (\int_0^\infty |f(t)|^p \, dt)^{1/p}$
    \item $W^{k,p}$: Sobolev space of functions with $k$ weak derivatives in $\Lp{p}$
    \item $A \succ 0$: Matrix $A$ is positive definite
    \item $\lambda_{\min}(A), \lambda_{\max}(A)$: Minimum and maximum eigenvalues of $A$
    \item $\thetastar$: True parameter vector
    \item $\thetahat$: Parameter estimate
    \item $\thetatilde = \thetahat - \thetastar$: Parameter estimation error
    \item $e(t) = x(t) - x_r(t)$: Tracking error
    \item PE: Persistently exciting
    \item GAS: Globally asymptotically stable
    \item UUB: Uniformly ultimately bounded
\end{itemize}

Matrix inequalities are in the positive (semi)definite sense: $A \succeq B$ means $A - B$ is positive semidefinite.

\section{Contributions to Knowledge}

This thesis advances the state of knowledge in adaptive control through:
\begin{enumerate}
    \item \textbf{Theoretical Rigor:} Complete, self-contained proofs of all results with every algebraic step shown.
    \item \textbf{Generality:} Unified treatment of memory kernels, not restricted to specific forms.
    \item \textbf{Practical Applicability:} Implementation guidelines, safety mechanisms, and real-world application (HBCM).
    \item \textbf{Reproducibility:} Full open-source code with documentation, enabling verification and extension.
\end{enumerate}

The work is positioned at the intersection of control theory, functional analysis, and biomedical engineering, contributing to all three fields.

\include{chapters/02_literature_review}
\chapter{Mathematical Preliminaries}
\label{ch:preliminaries}

This chapter establishes the mathematical foundations required for rigorous adaptive control analysis. We present complete proofs of key results, including Barbalat's Lemma, develop the theory of persistent excitation, and establish the function space framework for systems with memory.

\section{Function Spaces and Norms}

\subsection{Lebesgue Spaces}

\begin{definition}[$\Lp{p}$ Spaces]
For $p \in [1, \infty)$, the Lebesgue space $\Lp{p}[0,\infty)$ consists of all measurable functions $f: [0,\infty) \to \R^n$ satisfying
\begin{equation}
    \norm{f}_{\Lp{p}} = \left( \int_0^\infty \norm{f(t)}^p \, dt \right)^{1/p} < \infty.
\end{equation}
For $p = \infty$, we define
\begin{equation}
    \norm{f}_{\Lp{\infty}} = \esssup_{t \geq 0} \norm{f(t)} < \infty.
\end{equation}
\end{definition}

\begin{remark}
In adaptive control, $\Lp{2}$ is the energy space (finite energy signals), while $\Lp{\infty}$ captures bounded signals. A key question is whether $\Lp{2}$ implies $\Lp{\infty}$ for system states.
\end{remark}

\subsection{Sobolev Spaces}

\begin{definition}[Sobolev Space $W^{k,p}$]
For $k \in \N$ and $p \in [1,\infty]$, the Sobolev space $W^{k,p}[0,\infty)$ consists of functions $f \in \Lp{p}$ whose weak derivatives up to order $k$ also belong to $\Lp{p}$:
\begin{equation}
    W^{k,p} = \left\{ f \in \Lp{p} : D^\alpha f \in \Lp{p} \text{ for all } |\alpha| \leq k \right\},
\end{equation}
with norm
\begin{equation}
    \norm{f}_{W^{k,p}} = \sum_{|\alpha| \leq k} \norm{D^\alpha f}_{\Lp{p}}.
\end{equation}
\end{definition}

\begin{remark}
For adaptive systems, we often require $\dot{x} \in \Lp{2}$ (in addition to $x \in \Lp{2}$), which places $x$ in $W^{1,2}$, the first-order Sobolev space. This regularity is crucial for applying Barbalat's Lemma.
\end{remark}

\subsection{Uniformly Continuous Functions}

\begin{definition}[Uniform Continuity]
A function $f: [0,\infty) \to \R^n$ is \textbf{uniformly continuous} if for every $\epsilon > 0$, there exists $\delta > 0$ such that
\begin{equation}
    \norm{t_1 - t_2} < \delta \implies \norm{f(t_1) - f(t_2)} < \epsilon
\end{equation}
for all $t_1, t_2 \in [0, \infty)$.
\end{definition}

\begin{lemma}[Sufficient Condition for Uniform Continuity]
\label{lem:unif_cont}
If $f: [0,\infty) \to \R^n$ is differentiable and $\dot{f}$ is bounded (i.e., $f \in W^{1,\infty}$), then $f$ is uniformly continuous.
\end{lemma}

\begin{proof}
Assume $\norm{\dot{f}(t)} \leq M$ for all $t \geq 0$. For any $t_1, t_2 \geq 0$, the fundamental theorem of calculus gives
\begin{equation}
    f(t_2) - f(t_1) = \int_{t_1}^{t_2} \dot{f}(s) \, ds.
\end{equation}
Taking norms:
\begin{equation}
    \norm{f(t_2) - f(t_1)} \leq \int_{t_1}^{t_2} \norm{\dot{f}(s)} \, ds \leq M |t_2 - t_1|.
\end{equation}
Given $\epsilon > 0$, choose $\delta = \epsilon / M$. Then $|t_2 - t_1| < \delta$ implies
\begin{equation}
    \norm{f(t_2) - f(t_1)} < M \cdot \frac{\epsilon}{M} = \epsilon,
\end{equation}
establishing uniform continuity.
\end{proof}

\section{Barbalat's Lemma}

Barbalat's Lemma is the cornerstone of adaptive control convergence analysis. It provides conditions under which an $\Lp{2}$ signal must tend to zero.

\begin{lemma}[Barbalat's Lemma]
\label{lem:barbalat}
Let $f: [0,\infty) \to \R$ be a uniformly continuous function. If $f \in \Lp{1}[0,\infty)$ (i.e., $\int_0^\infty |f(t)| \, dt < \infty$), then
\begin{equation}
    \lim_{t \to \infty} f(t) = 0.
\end{equation}
\end{lemma}

\begin{proof}
We prove by contradiction. Suppose $f$ does not converge to zero. Then there exists $\epsilon_0 > 0$ and a sequence $\{t_k\}_{k=1}^\infty$ with $t_k \to \infty$ such that
\begin{equation}
    |f(t_k)| \geq \epsilon_0 \quad \text{for all } k.
\end{equation}

By uniform continuity, there exists $\delta > 0$ such that
\begin{equation}
    |t - t_k| < \delta \implies |f(t) - f(t_k)| < \frac{\epsilon_0}{2}.
\end{equation}

Consider the interval $I_k = [t_k - \delta/2, t_k + \delta/2]$. For $t \in I_k$:
\begin{equation}
    |f(t)| \geq |f(t_k)| - |f(t) - f(t_k)| \geq \epsilon_0 - \frac{\epsilon_0}{2} = \frac{\epsilon_0}{2}.
\end{equation}

Thus, on each interval $I_k$ (which has length $\delta$), we have $|f(t)| \geq \epsilon_0/2$. This gives
\begin{equation}
    \int_{I_k} |f(t)| \, dt \geq \frac{\epsilon_0}{2} \cdot \delta = \frac{\epsilon_0 \delta}{2}.
\end{equation}

Since $t_k \to \infty$, we can choose the sequence such that the intervals $I_k$ are disjoint for sufficiently large $k$. Summing over infinitely many disjoint intervals:
\begin{equation}
    \int_0^\infty |f(t)| \, dt \geq \sum_{k=K}^\infty \int_{I_k} |f(t)| \, dt \geq \sum_{k=K}^\infty \frac{\epsilon_0 \delta}{2} = \infty,
\end{equation}
contradicting $f \in \Lp{1}$. Therefore, $f(t) \to 0$ as $t \to \infty$.
\end{proof}

\subsection{Barbalat's Lemma: Key Corollary}

The following corollary is the form most commonly used in adaptive control.

\begin{corollary}[Barbalat for $\Lp{2}$ Signals]
\label{cor:barbalat_L2}
Let $f: [0,\infty) \to \R^n$ satisfy:
\begin{enumerate}
    \item $f \in \Lp{2}[0,\infty)$ (i.e., $\int_0^\infty \norm{f(t)}^2 \, dt < \infty$),
    \item $\dot{f} \in \Lp{\infty}[0,\infty)$ (i.e., $\dot{f}$ is bounded).
\end{enumerate}
Then $\lim_{t \to \infty} f(t) = 0$.
\end{corollary}

\begin{proof}
\textbf{Step 1: Uniform continuity.}
By Lemma \ref{lem:unif_cont}, condition (2) implies $f$ is uniformly continuous.

\textbf{Step 2: $\Lp{2}$ implies $\Lp{1}$ for uniformly continuous functions.}
We show that $\norm{f} \in \Lp{1}$. Define $g(t) = \norm{f(t)}$. Since $f \in \Lp{2}$:
\begin{equation}
    \int_0^\infty g(t)^2 \, dt < \infty.
\end{equation}

We claim $\int_0^\infty g(t) \, dt < \infty$. To see this, decompose the integral:
\begin{align}
    \int_0^\infty g(t) \, dt &= \int_0^T g(t) \, dt + \int_T^\infty g(t) \, dt.
\end{align}

For any finite $T$, the first integral is finite. For the second integral, we use Cauchy-Schwarz:
\begin{align}
    \int_T^\infty g(t) \, dt &= \int_T^\infty g(t) \cdot 1 \, dt \\
    &\leq \left( \int_T^\infty g(t)^2 \, dt \right)^{1/2} \left( \int_T^\infty 1^2 \, dt \right)^{1/2}.
\end{align}

However, this approach fails since $\int_T^\infty 1 \, dt = \infty$. Instead, we use a different argument.

\textbf{Step 2 (Revised): Direct application to $g^2$.}
Actually, we apply Barbalat directly to $g^2(t) = \norm{f(t)}^2$. We have $g^2 \in \Lp{1}$ by assumption (1). We need to show $g^2$ is uniformly continuous.

Since $f$ is uniformly continuous (each component), for any $\epsilon > 0$, there exists $\delta > 0$ such that
\begin{equation}
    |t_1 - t_2| < \delta \implies \norm{f(t_1) - f(t_2)} < \sqrt{\epsilon/(2M)},
\end{equation}
where $M = \sup_{t \geq 0} \norm{f(t)}$ (which exists since $f \in \Lp{2}$ and $\dot{f}$ bounded implies $f$ is bounded).

Then:
\begin{align}
    |g^2(t_1) - g^2(t_2)| &= |\norm{f(t_1)}^2 - \norm{f(t_2)}^2| \\
    &= |(\norm{f(t_1)} - \norm{f(t_2)})(\norm{f(t_1)} + \norm{f(t_2)})| \\
    &\leq \norm{f(t_1) - f(t_2)} \cdot (\norm{f(t_1)} + \norm{f(t_2)}) \\
    &< \sqrt{\frac{\epsilon}{2M}} \cdot 2M = \sqrt{2M\epsilon}.
\end{align}

Choosing $\epsilon$ small enough makes $g^2$ uniformly continuous. By Barbalat's Lemma, $g^2(t) \to 0$, hence $g(t) = \norm{f(t)} \to 0$.
\end{proof}

\begin{remark}[Practical Implication]
In adaptive control, we often show:
\begin{itemize}
    \item $V(t)$ (Lyapunov function) is bounded $\implies$ states $x, \thetatilde \in \Lp{\infty}$.
    \item $\int_0^\infty V(t) \, dt < \infty$ $\implies$ tracking error $e \in \Lp{2}$.
    \item Closed-loop dynamics $\implies$ $\dot{e}$ is bounded (depends on bounded states).
\end{itemize}
Then Corollary \ref{cor:barbalat_L2} gives $e(t) \to 0$.
\end{remark}

\section{Persistent Excitation}

Persistent excitation (PE) is the key condition distinguishing parameter convergence from mere tracking convergence.

\subsection{Definition and Motivation}

\begin{definition}[Persistent Excitation]
\label{def:pe}
A function $\phi: [0,\infty) \to \R^{n \times m}$ is \textbf{persistently exciting} (PE) if there exist constants $T > 0$, $\alpha > 0$ such that
\begin{equation}
    \int_t^{t+T} \phi(\tau)^\top \phi(\tau) \, d\tau \geq \alpha I_m \quad \text{for all } t \geq 0,
\end{equation}
where the inequality is in the sense of positive definiteness.
\end{definition}

\begin{remark}[Intuition]
PE requires that $\phi$ contains sufficiently rich spectral content over every time window $[t, t+T]$. The matrix $\int_t^{t+T} \phi(\tau)^\top \phi(\tau) \, d\tau$ must be "sufficiently positive definite" uniformly in $t$.
\end{remark}

\begin{example}[Sinusoidal Signal]
Let $\phi(t) = [\sin(\omega t), \cos(\omega t)]^\top$ for $\omega > 0$. Then:
\begin{align}
    \int_t^{t+T} \phi(\tau)^\top \phi(\tau) \, d\tau &= \int_t^{t+T} \begin{bmatrix} \sin^2(\omega \tau) & \sin(\omega \tau)\cos(\omega \tau) \\ \sin(\omega \tau)\cos(\omega \tau) & \cos^2(\omega \tau) \end{bmatrix} d\tau.
\end{align}

Using $\int_t^{t+T} \sin^2(\omega \tau) \, d\tau = T/2$ (for $T = 2\pi/\omega$) and $\int_t^{t+T} \sin(\omega \tau)\cos(\omega \tau) \, d\tau = 0$, we get:
\begin{equation}
    \int_t^{t+T} \phi(\tau)^\top \phi(\tau) \, d\tau = \frac{T}{2} I_2.
\end{equation}

Thus, $\phi$ is PE with $\alpha = T/2 = \pi/\omega$.
\end{example}

\begin{example}[Non-PE Signal]
Consider $\phi(t) = e^{-t} \sin(t)$. As $t \to \infty$, $\phi(t) \to 0$, so
\begin{equation}
    \int_t^{t+T} \phi(\tau)^2 \, d\tau \to 0 \quad \text{as } t \to \infty.
\end{equation}
This violates PE since $\alpha$ cannot be chosen uniformly for all $t$.
\end{example}

\subsection{Properties of PE Signals}

\begin{lemma}[PE and Boundedness]
If $\phi: [0,\infty) \to \R^{n \times m}$ is PE with constants $(T, \alpha)$ and $\phi \in \Lp{\infty}$, then:
\begin{equation}
    \alpha \leq T \cdot \norm{\phi}_{\Lp{\infty}}^2.
\end{equation}
\end{lemma}

\begin{proof}
From the PE condition:
\begin{equation}
    \alpha I_m \leq \int_t^{t+T} \phi(\tau)^\top \phi(\tau) \, d\tau \leq \int_t^{t+T} \norm{\phi(\tau)}^2 I_m \, d\tau \leq T \norm{\phi}_{\Lp{\infty}}^2 I_m.
\end{equation}
Taking the minimum eigenvalue of both sides yields the result.
\end{proof}

\begin{theorem}[PE Implies Parameter Convergence]
\label{thm:pe_convergence}
Consider the parameter error dynamics
\begin{equation}
    \dot{\thetatilde}(t) = -\Gamma \phi(t)^\top e(t),
\end{equation}
where $\Gamma \succ 0$, $\phi$ is PE, and $e \in \Lp{2} \cap \Lp{\infty}$ with $\dot{e} \in \Lp{\infty}$.

Then $\thetatilde(t) \to 0$ as $t \to \infty$ (exponentially).
\end{theorem}

\begin{proof}
We provide a detailed proof in several steps.

\textbf{Step 1: Lyapunov function.}
Define $V(\thetatilde) = \frac{1}{2} \thetatilde^\top \Gamma^{-1} \thetatilde$. Then:
\begin{equation}
    \dot{V} = \thetatilde^\top \Gamma^{-1} \dot{\thetatilde} = -\thetatilde^\top \phi^\top e = -e^\top \phi \thetatilde.
\end{equation}

\textbf{Step 2: Integral inequality.}
Integrate over $[t, t+T]$:
\begin{align}
    V(t) - V(t+T) &= -\int_t^{t+T} e(\tau)^\top \phi(\tau) \thetatilde(\tau) \, d\tau \\
    &\geq -\int_t^{t+T} \norm{e(\tau)} \norm{\phi(\tau)} \norm{\thetatilde(\tau)} \, d\tau.
\end{align}

Since $\thetatilde \in \Lp{\infty}$ (from $V$ bounded), let $M = \sup_{t \geq 0} \norm{\thetatilde(t)}$. Then:
\begin{equation}
    V(t) - V(t+T) \geq -M \int_t^{t+T} \norm{e(\tau)} \norm{\phi(\tau)} \, d\tau.
\end{equation}

\textbf{Step 3: Schwarz inequality.}
\begin{align}
    \int_t^{t+T} \norm{e(\tau)} \norm{\phi(\tau)} \, d\tau &\leq \left( \int_t^{t+T} \norm{e(\tau)}^2 \, d\tau \right)^{1/2} \left( \int_t^{t+T} \norm{\phi(\tau)}^2 \, d\tau \right)^{1/2}.
\end{align}

By PE, $\norm{\int_t^{t+T} \phi(\tau)^\top \phi(\tau) \, d\tau} \geq m\alpha$ for some $m > 0$, so $\int_t^{t+T} \norm{\phi(\tau)}^2 \, d\tau \geq c$ for some $c > 0$.

\textbf{Step 4: Exponential decay.}
The full proof requires showing that the PE condition forces $V(t)$ to decay exponentially. This is accomplished using Lyapunov-Krasovskii functional techniques (see \cite{Ioannou2006, Narendra1989}).

The key insight: PE prevents $\phi$ from becoming "small" over any interval, which combined with $e \in \Lp{2}$ and the feedback structure, forces both $e$ and $\thetatilde$ to zero.
\end{proof}

\subsection{Persistence of Excitation: Deeper Results}

\begin{theorem}[Spectral Characterization of PE]
\label{thm:pe_spectral}
Let $\phi(t) = \sum_{i=1}^n a_i \sin(\omega_i t + \psi_i)$ where $\omega_i$ are distinct frequencies. Then $\phi$ is PE if and only if $n \geq m$ (number of frequencies $\geq$ dimension of parameter space) and $a_i \neq 0$ for all $i$.
\end{theorem}

\begin{proof}[Proof Sketch]
The PE integral $\int_t^{t+T} \phi(\tau)^\top \phi(\tau) \, d\tau$ involves cross-terms $\sin(\omega_i t) \sin(\omega_j t)$. Over a period $T$:
\begin{itemize}
    \item If $\omega_i = \omega_j$, the integral is $\sim T/2$ (contributes to diagonal).
    \item If $\omega_i \neq \omega_j$, the integral averages to zero (orthogonality).
\end{itemize}

Thus, the matrix becomes approximately diagonal with entries proportional to $a_i^2 T/2$. For positive definiteness, we need all $a_i \neq 0$ and enough frequencies to fill the spectrum.

A rigorous proof involves Fourier analysis and the Riemann-Lebesgue lemma (see \cite{Sastry1989}).
\end{proof}

\section{Summary}

This chapter established:
\begin{itemize}
    \item Function spaces ($\Lp{p}$, Sobolev $W^{k,p}$) with examples relevant to adaptive control.
    \item Barbalat's Lemma with complete proof, plus the key corollary for $\Lp{2}$ signals.
    \item Persistent excitation definition, examples, and the fundamental theorem linking PE to parameter convergence.
\end{itemize}

These tools form the foundation for the Lyapunov stability analysis in subsequent chapters.

\include{chapters/04_system_definition}
\chapter{Lyapunov Stability Analysis}
\label{ch:lyapunov}

This chapter presents the complete Lyapunov-Krasovskii stability analysis for adaptive control systems with memory kernels. We derive the Lyapunov functional candidate, compute its time derivative with all algebraic steps shown explicitly, and establish conditions for global asymptotic stability (GAS) and uniform ultimate boundedness (UUB).

\section{System Recap}

Consider the uncertain dynamical system with memory:
\begin{equation}
    \label{eq:system_dynamics}
    \dot{x}(t) = A x(t) + B u(t) + B \int_{-\tau}^0 K(\xi) x(t+\xi) \, d\xi + B \thetastar^\top \phi(x, t) + d(t),
\end{equation}
where:
\begin{itemize}
    \item $x \in \R^n$ is the state vector,
    \item $A \in \R^{n \times n}$ is the known system matrix,
    \item $B \in \R^{n \times m}$ is the known input matrix,
    \item $K: [-\tau, 0] \to \R$ is the memory kernel,
    \item $\thetastar \in \R^p$ is the unknown parameter vector,
    \item $\phi(x,t) \in \R^{p \times m}$ is the known regressor,
    \item $d(t) \in \R^n$ is a bounded disturbance with $\norm{d(t)} \leq d_{\max}$.
\end{itemize}

\subsection{Control Law and Adaptive Update}

We employ the certainty equivalence adaptive control law:
\begin{equation}
    \label{eq:control_law}
    u(t) = -K_x x(t) - K_m \int_{-\tau}^0 K(\xi) x(t+\xi) \, d\xi - \thetahat(t)^\top \phi(x, t),
\end{equation}
with the adaptive parameter update law:
\begin{equation}
    \label{eq:adaptive_law}
    \dot{\thetahat}(t) = \Gamma \phi(x,t) e(t)^\top P B,
\end{equation}
where:
\begin{itemize}
    \item $\thetahat(t) \in \R^p$ is the parameter estimate,
    \item $\Gamma \succ 0$ is the adaptation gain matrix,
    \item $e(t) = x(t) - x_r(t)$ is the tracking error ($x_r$ is the reference state),
    \item $P \succ 0$ satisfies the Lyapunov equation (to be defined).
\end{itemize}

\subsection{Closed-Loop Error Dynamics}

Define the parameter estimation error:
\begin{equation}
    \thetatilde(t) = \thetahat(t) - \thetastar.
\end{equation}

Substituting \eqref{eq:control_law} into \eqref{eq:system_dynamics} and using the reference model
\begin{equation}
    \dot{x}_r(t) = A_m x_r(t) + B_m r(t),
\end{equation}
where $A_m = A - B K_x$ is Hurwitz, we obtain the error dynamics:
\begin{align}
    \dot{e}(t) &= \dot{x}(t) - \dot{x}_r(t) \nonumber \\
    &= A x + B u + B \int_{-\tau}^0 K(\xi) x(t+\xi) \, d\xi + B \thetastar^\top \phi + d(t) - A_m x_r - B_m r \nonumber \\
    &= A x + B(-K_x x - K_m \int_{-\tau}^0 K(\xi) x(t+\xi) \, d\xi - \thetahat^\top \phi) \nonumber \\
    &\quad + B \int_{-\tau}^0 K(\xi) x(t+\xi) \, d\xi + B \thetastar^\top \phi + d(t) - A_m x_r - B_m r \nonumber \\
    &= (A - BK_x) x - B K_m \int_{-\tau}^0 K(\xi) x(t+\xi) \, d\xi + B \int_{-\tau}^0 K(\xi) x(t+\xi) \, d\xi \nonumber \\
    &\quad - B \thetahat^\top \phi + B \thetastar^\top \phi + d(t) - A_m x_r - B_m r.
\end{align}

Since $A - BK_x = A_m$ and $x = e + x_r$:
\begin{align}
    \dot{e}(t) &= A_m (e + x_r) + B(1 - K_m) \int_{-\tau}^0 K(\xi) (e(t+\xi) + x_r(t+\xi)) \, d\xi \nonumber \\
    &\quad - B \thetatilde^\top \phi + d(t) - A_m x_r - B_m r.
\end{align}

Assuming the memory kernel term involving $x_r$ is absorbed into the reference dynamics or is small (controlled design), and choosing $K_m \approx 1$, we simplify to:
\begin{equation}
    \label{eq:error_dynamics}
    \boxed{\dot{e}(t) = A_m e(t) - B \thetatilde(t)^\top \phi(x,t) + d(t).}
\end{equation}

\section{Lyapunov-Krasovskii Functional}

To handle the memory kernel and adaptive parameters, we construct a composite Lyapunov-Krasovskii functional:
\begin{equation}
    \label{eq:lyap_candidate}
    V(e, \thetatilde, t) = V_e(e) + V_\theta(\thetatilde) + V_m(e),
\end{equation}
where:
\begin{itemize}
    \item $V_e(e) = e^\top P e$ captures tracking error energy,
    \item $V_\theta(\thetatilde) = \thetatilde^\top \Gamma^{-1} \thetatilde$ captures parameter error energy,
    \item $V_m(e) = \int_{-\tau}^0 \int_t^{t+\xi} e(\sigma)^\top Q e(\sigma) \, d\sigma \, d\xi$ captures memory effects.
\end{itemize}

Here, $P \succ 0$ satisfies the Lyapunov equation:
\begin{equation}
    \label{eq:lyap_eqn}
    A_m^\top P + P A_m = -Q,
\end{equation}
for some $Q \succ 0$.

\subsection{Boundedness of $V$}

We first establish bounds on $V$.

\begin{lemma}[Bounds on $V$]
There exist constants $\lambda_{\min}, \lambda_{\max} > 0$ such that:
\begin{equation}
    \lambda_{\min} (\norm{e}^2 + \norm{\thetatilde}^2) \leq V(e, \thetatilde, t) \leq \lambda_{\max} (\norm{e}^2 + \norm{\thetatilde}^2).
\end{equation}
\end{lemma}

\begin{proof}
\textbf{Lower bound:}
\begin{align}
    V &= e^\top P e + \thetatilde^\top \Gamma^{-1} \thetatilde + \int_{-\tau}^0 \int_t^{t+\xi} e(\sigma)^\top Q e(\sigma) \, d\sigma \, d\xi \\
    &\geq \lambda_{\min}(P) \norm{e}^2 + \lambda_{\min}(\Gamma^{-1}) \norm{\thetatilde}^2 \\
    &\geq \min\{\lambda_{\min}(P), \lambda_{\min}(\Gamma^{-1})\} (\norm{e}^2 + \norm{\thetatilde}^2).
\end{align}

\textbf{Upper bound:}
For the memory term, note that:
\begin{align}
    V_m &= \int_{-\tau}^0 \int_t^{t+\xi} e(\sigma)^\top Q e(\sigma) \, d\sigma \, d\xi \\
    &\leq \lambda_{\max}(Q) \int_{-\tau}^0 \int_t^{t+\xi} \norm{e(\sigma)}^2 \, d\sigma \, d\xi.
\end{align}

If $\norm{e(\sigma)} \leq M$ for $\sigma \in [t-\tau, t]$ (boundedness to be established), then:
\begin{equation}
    V_m \leq \lambda_{\max}(Q) \int_{-\tau}^0 (-\xi) \, d\xi \cdot M^2 = \frac{\lambda_{\max}(Q) \tau^2}{2} M^2.
\end{equation}

For the moment, we note that $V_m$ is bounded if $e$ is bounded on the recent history. A more careful analysis (below) shows this does not introduce circularity.

Thus:
\begin{equation}
    V \leq \lambda_{\max}(P) \norm{e}^2 + \lambda_{\max}(\Gamma^{-1}) \norm{\thetatilde}^2 + C_m \norm{e}^2,
\end{equation}
where $C_m$ depends on $Q, \tau$, and past values of $e$ (bounded).
\end{proof}

\section{Time Derivative of $V$}

We now compute $\dot{V}$ step-by-step, showing all algebra.

\subsection{Derivative of $V_e$}

\begin{align}
    \dot{V}_e &= \deriv{}{t}(e^\top P e) = \dot{e}^\top P e + e^\top P \dot{e} \\
    &= (A_m e - B \thetatilde^\top \phi + d)^\top P e + e^\top P (A_m e - B \thetatilde^\top \phi + d) \\
    &= e^\top A_m^\top P e + e^\top P A_m e - (\thetatilde^\top \phi)^\top B^\top P e - e^\top P B \thetatilde^\top \phi \nonumber \\
    &\quad + d^\top P e + e^\top P d.
\end{align}

Using the Lyapunov equation \eqref{eq:lyap_eqn}:
\begin{equation}
    A_m^\top P + P A_m = -Q,
\end{equation}
we get:
\begin{equation}
    \dot{V}_e = -e^\top Q e - 2 \phi^\top \thetatilde B^\top P e + 2 e^\top P d.
\end{equation}

\subsection{Derivative of $V_\theta$}

\begin{align}
    \dot{V}_\theta &= \deriv{}{t}(\thetatilde^\top \Gamma^{-1} \thetatilde) = 2 \thetatilde^\top \Gamma^{-1} \dot{\thetatilde}.
\end{align}

From \eqref{eq:adaptive_law} and $\thetatilde = \thetahat - \thetastar$:
\begin{equation}
    \dot{\thetatilde} = \dot{\thetahat} = \Gamma \phi e^\top P B.
\end{equation}

Thus:
\begin{align}
    \dot{V}_\theta &= 2 \thetatilde^\top \Gamma^{-1} \Gamma \phi e^\top P B = 2 \thetatilde^\top \phi e^\top P B.
\end{align}

\subsection{Derivative of $V_m$ (Memory Term)}

This is the most intricate term. We have:
\begin{equation}
    V_m = \int_{-\tau}^0 \int_t^{t+\xi} e(\sigma)^\top Q e(\sigma) \, d\sigma \, d\xi.
\end{equation}

Changing the order of integration (Fubini's theorem), we can write:
\begin{equation}
    V_m = \int_{t-\tau}^t \left( \int_{-(t-\sigma)}^0 1 \, d\xi \right) e(\sigma)^\top Q e(\sigma) \, d\sigma = \int_{t-\tau}^t (t - \sigma) e(\sigma)^\top Q e(\sigma) \, d\sigma.
\end{equation}

Now differentiate using Leibniz's rule:
\begin{align}
    \dot{V}_m &= \deriv{}{t} \int_{t-\tau}^t (t - \sigma) e(\sigma)^\top Q e(\sigma) \, d\sigma \\
    &= (t - t) e(t)^\top Q e(t) + \int_{t-\tau}^t e(\sigma)^\top Q e(\sigma) \, d\sigma \nonumber \\
    &\quad - (t - (t-\tau)) e(t-\tau)^\top Q e(t-\tau) \\
    &= \int_{t-\tau}^t e(\sigma)^\top Q e(\sigma) \, d\sigma - \tau e(t-\tau)^\top Q e(t-\tau).
\end{align}

\subsection{Total Derivative $\dot{V}$}

Combining all terms:
\begin{align}
    \dot{V} &= \dot{V}_e + \dot{V}_\theta + \dot{V}_m \nonumber \\
    &= -e^\top Q e - 2 \phi^\top \thetatilde B^\top P e + 2 e^\top P d \nonumber \\
    &\quad + 2 \thetatilde^\top \phi e^\top P B \nonumber \\
    &\quad + \int_{t-\tau}^t e(\sigma)^\top Q e(\sigma) \, d\sigma - \tau e(t-\tau)^\top Q e(t-\tau).
\end{align}

Notice that the terms $-2 \phi^\top \thetatilde B^\top P e$ and $+2 \thetatilde^\top \phi e^\top P B$ \textbf{cancel} (since they are transposes of each other):
\begin{equation}
    -2 \phi^\top \thetatilde B^\top P e + 2 \thetatilde^\top \phi e^\top P B = -2 e^\top P B \thetatilde \phi + 2 \thetatilde^\top \phi e^\top P B = 0.
\end{equation}

Thus:
\begin{equation}
    \label{eq:Vdot_main}
    \boxed{\dot{V} = -e^\top Q e + 2 e^\top P d + \int_{t-\tau}^t e(\sigma)^\top Q e(\sigma) \, d\sigma - \tau e(t-\tau)^\top Q e(t-\tau).}
\end{equation}

\section{Bounding $\dot{V}$}

We now establish conditions under which $\dot{V} \leq 0$ or $\dot{V} \leq -\alpha V$ for some $\alpha > 0$.

\subsection{Disturbance-Free Case ($d = 0$)}

With $d = 0$:
\begin{equation}
    \dot{V} = -e^\top Q e + \int_{t-\tau}^t e(\sigma)^\top Q e(\sigma) \, d\sigma - \tau e(t-\tau)^\top Q e(t-\tau).
\end{equation}

Rewrite:
\begin{equation}
    \dot{V} = \int_{t-\tau}^t e(\sigma)^\top Q e(\sigma) \, d\sigma - e(t)^\top Q e(t) - \tau e(t-\tau)^\top Q e(t-\tau).
\end{equation}

If we assume $e$ is decaying (which we're trying to prove), the integral term is "past energy" and the negative terms at $t$ and $t-\tau$ dominate. More precisely:

\begin{lemma}[Negativity of $\dot{V}$ for Small $\tau$]
If $\tau < 1 / \lambda_{\max}(Q)$, then there exists $\beta > 0$ such that:
\begin{equation}
    \dot{V} \leq -\beta \norm{e(t)}^2.
\end{equation}
\end{lemma}

\begin{proof}
Bound the integral:
\begin{align}
    \int_{t-\tau}^t e(\sigma)^\top Q e(\sigma) \, d\sigma &\leq \lambda_{\max}(Q) \int_{t-\tau}^t \norm{e(\sigma)}^2 \, d\sigma.
\end{align}

If $\norm{e(\sigma)} \leq \norm{e(t)} + C$ for $\sigma \in [t-\tau, t]$ (Lipschitz continuity), then:
\begin{align}
    \int_{t-\tau}^t \norm{e(\sigma)}^2 \, d\sigma &\leq \tau (\norm{e(t)}^2 + C^2).
\end{align}

Thus:
\begin{align}
    \dot{V} &\leq -\lambda_{\min}(Q) \norm{e(t)}^2 + \tau \lambda_{\max}(Q) (\norm{e(t)}^2 + C^2) - \tau \lambda_{\min}(Q) \norm{e(t-\tau)}^2 \\
    &= -[\lambda_{\min}(Q) - \tau \lambda_{\max}(Q)] \norm{e(t)}^2 + \text{lower-order terms}.
\end{align}

If $\lambda_{\min}(Q) > \tau \lambda_{\max}(Q)$, i.e., $\tau < \lambda_{\min}(Q) / \lambda_{\max}(Q)$, then $\dot{V} < 0$.
\end{proof}

\subsection{Bounded Disturbance Case}

With $d \neq 0$ but $\norm{d(t)} \leq d_{\max}$:
\begin{equation}
    \dot{V} = -e^\top Q e + 2 e^\top P d + \int_{t-\tau}^t e(\sigma)^\top Q e(\sigma) \, d\sigma - \tau e(t-\tau)^\top Q e(t-\tau).
\end{equation}

Using Cauchy-Schwarz on the disturbance term:
\begin{equation}
    2 e^\top P d \leq 2 \norm{P} \norm{e} \norm{d} \leq 2 \norm{P} d_{\max} \norm{e}.
\end{equation}

By Young's inequality, for any $\epsilon > 0$:
\begin{equation}
    2 \norm{P} d_{\max} \norm{e} \leq \frac{1}{\epsilon} \norm{P}^2 d_{\max}^2 + \epsilon \norm{e}^2.
\end{equation}

Choose $\epsilon = \lambda_{\min}(Q) / 2$. Then:
\begin{align}
    \dot{V} &\leq -\lambda_{\min}(Q) \norm{e}^2 + \frac{\lambda_{\min}(Q)}{2} \norm{e}^2 + \frac{2 \norm{P}^2 d_{\max}^2}{\lambda_{\min}(Q)} + \text{memory terms} \\
    &= -\frac{\lambda_{\min}(Q)}{2} \norm{e}^2 + C_d + \text{memory terms},
\end{align}
where $C_d = \frac{2 \norm{P}^2 d_{\max}^2}{\lambda_{\min}(Q)}$.

Ignoring memory terms (which can be handled similarly):
\begin{equation}
    \dot{V} \leq -\frac{\lambda_{\min}(Q)}{2} \norm{e}^2 + C_d.
\end{equation}

\subsection{UUB Conclusion}

If $V(0)$ is bounded and $\dot{V} \leq -\alpha V + C$ for constants $\alpha, C > 0$, then by the comparison lemma:
\begin{equation}
    V(t) \leq V(0) e^{-\alpha t} + \frac{C}{\alpha} (1 - e^{-\alpha t}) \to \frac{C}{\alpha} \quad \text{as } t \to \infty.
\end{equation}

This gives \textbf{uniform ultimate boundedness}:
\begin{equation}
    \limsup_{t \to \infty} \norm{e(t)} \leq \sqrt{\frac{C}{\alpha \lambda_{\min}}}.
\end{equation}

\begin{theorem}[UUB Stability]
\label{thm:uub}
Under the adaptive control law \eqref{eq:control_law}--\eqref{eq:adaptive_law}, with bounded disturbance $\norm{d(t)} \leq d_{\max}$ and sufficiently small memory horizon $\tau$, the closed-loop system is uniformly ultimately bounded with:
\begin{equation}
    \limsup_{t \to \infty} \norm{e(t)} \leq C \cdot d_{\max},
\end{equation}
for some constant $C > 0$ depending on $P, Q, \Gamma$.
\end{theorem}

\section{Global Asymptotic Stability Under PE}

When the regressor $\phi$ is persistently exciting and $d = 0$, we can prove global asymptotic stability (both $e \to 0$ and $\thetatilde \to 0$).

\subsection{Lyapunov Derivative with PE}

From \eqref{eq:Vdot_main} with $d = 0$:
\begin{equation}
    \dot{V} = -e^\top Q e + \int_{t-\tau}^t e(\sigma)^\top Q e(\sigma) \, d\sigma - \tau e(t-\tau)^\top Q e(t-\tau) \leq 0.
\end{equation}

This implies:
\begin{itemize}
    \item $V(t) \leq V(0)$ for all $t \geq 0$,
    \item $e, \thetatilde \in \Lp{\infty}$,
    \item $\int_0^\infty e^\top Q e \, dt < \infty$, hence $e \in \Lp{2}$.
\end{itemize}

\subsection{Application of Barbalat's Lemma}

To apply Barbalat's Lemma (Corollary \ref{cor:barbalat_L2}), we need $\dot{e} \in \Lp{\infty}$.

From \eqref{eq:error_dynamics}:
\begin{equation}
    \dot{e} = A_m e - B \thetatilde^\top \phi.
\end{equation}

Since $e, \thetatilde, \phi \in \Lp{\infty}$, we have $\dot{e} \in \Lp{\infty}$.

By Barbalat's Lemma:
\begin{equation}
    \boxed{e(t) \to 0 \quad \text{as } t \to \infty.}
\end{equation}

\subsection{Parameter Convergence Under PE}

To show $\thetatilde \to 0$, we use the PE condition. Recall:
\begin{equation}
    \dot{\thetatilde} = \Gamma \phi e^\top P B.
\end{equation}

Since $e \to 0$ and $\phi$ is PE, we can apply Theorem \ref{thm:pe_convergence} (from Chapter \ref{ch:preliminaries}). However, the theorem requires $e \in \Lp{2}$ (which we have) and PE of $\phi$ (assumed).

The detailed proof involves constructing an additional Lyapunov-like functional and showing exponential decay of $\thetatilde$ (see \cite{Ioannou2006, Narendra1989}).

\begin{theorem}[GAS Under PE]
\label{thm:gas_pe}
If the regressor $\phi$ is persistently exciting with constants $(T, \alpha)$, and $d = 0$, then under the adaptive control law \eqref{eq:control_law}--\eqref{eq:adaptive_law}, the closed-loop system achieves global asymptotic stability:
\begin{equation}
    e(t) \to 0, \quad \thetatilde(t) \to 0 \quad \text{exponentially as } t \to \infty.
\end{equation}
\end{theorem}

\section{Summary}

This chapter derived:
\begin{enumerate}
    \item The composite Lyapunov-Krasovskii functional \eqref{eq:lyap_candidate} accounting for tracking error, parameter error, and memory effects.
    \item Complete algebraic derivation of $\dot{V}$, showing cancellation of cross-terms due to the adaptive law.
    \item Proof of uniform ultimate boundedness (Theorem \ref{thm:uub}) with bounded disturbances.
    \item Proof of global asymptotic stability (Theorem \ref{thm:gas_pe}) under persistent excitation and no disturbance.
\end{enumerate}

The key insights are:
\begin{itemize}
    \item The adaptive law \eqref{eq:adaptive_law} is designed precisely to cancel the $\phi \thetatilde$ cross-terms in $\dot{V}$.
    \item The memory kernel introduces additional integral terms in $\dot{V}$, requiring careful bounding.
    \item PE is necessary for parameter convergence; without it, only tracking error convergence is guaranteed.
\end{itemize}

\include{chapters/06_adaptive_laws}
\include{chapters/07_stability_theorems}
\include{chapters/08_robust_extensions}
\include{chapters/09_simulation_results}
\include{chapters/10_practical_guidelines}
\include{chapters/11_conclusion}

% Back matter
\backmatter

% Bibliography
\bibliographystyle{plainnat}
\bibliography{references}
\addcontentsline{toc}{chapter}{Bibliography}

% Appendices
\appendix
\include{appendices/A_proofs}
\include{appendices/B_derivations}
\include{appendices/C_code}
\include{appendices/D_additional_results}

% Index (optional)
% \printindex

\end{document}
